\begin{quote}
\textbf{Maturity: {[}Practiced{]}, strongly.} This is a practitioner
essay, not a controlled study. What has changed since February 2026 is
that the workflow it describes was subsequently codified into the
project's standing engineering framework and used to build the rest of
the series' software. This revision reports that maturation honestly,
and keeps the essay's original candor: the field observations are
observations, not statistically significant results.
\end{quote}

\subsection{Abstract}\label{abstract}

AI coding agents can generate functionally correct code but lack the
engineering discipline that prevents scope drift, architectural
violations, and silent regressions over extended development sessions.
This paper argues that Behavior-Driven Development (BDD), specifically
the Given/When/Then specification language (Gherkin) combined with the
red-green-refactor cycle, provides the missing structural discipline for
agentic engineering. The argument is grounded in practitioner
experience: the author ran three concurrent agentic builds using
BDD-first methodology across multiple agent platforms and observed
consistent patterns of improvement when agents were constrained by
executable specifications rather than free-form instructions. We
formalize these observations as a methodology and connect them to the
Citrate Network's Mentorship Protocol (Paper III): the test suite is the
mentor, the coding agent is the mentee, and passing tests are successful
knowledge transfer. Unlike the February 2026 draft, we can now report
the sequel: this methodology was codified into the project's
\textbf{Agentile} engineering framework, a set of non-negotiable rules
and workflows, and used to build the software behind Papers I--III and
X--XI. This remains a practitioner report; we describe the experimental
methodology that would validate the observations, and we have not run
it.

\textbf{Keywords:} behavior-driven development, BDD, Gherkin, agentic
systems, test-driven development, AI coding agents, red-green-refactor,
specification language, software engineering methodology

\subsection{1. Introduction}\label{introduction}

The AI coding-agent shift inverted a classical assumption: that the
bottleneck in development is writing code. Agents generate hundreds of
lines of plausible code per minute. The new bottleneck is
\textbf{specification}, telling the agent what to build precisely enough
that the output matches intent, integrates with the existing
architecture, and does not silently break what already worked.

This paper's specific claim: Behavior-Driven Development, as formalized
by Dan North {[}1{]} and implemented through Gherkin {[}2{]}, is a
natural and effective discipline for constraining agentic code
generation. The Given/When/Then template forces the developer to specify
observable behavior before the agent writes implementation. The
red-green-refactor cycle, write a failing test, make it pass, clean up,
provides incremental checkpoints that prevent scope drift. And the
executable specification (\texttt{.feature} files) is living
documentation that persists across agent sessions, addressing the
context-window problem that plagues long-running builds.

\textbf{Practitioner basis.} The observations come from three concurrent
projects run with BDD-first methodology and AI coding agents between
October 2025 and February 2026: a HIPAA-compliant healthcare chatbot
(IKWE.ai, a separate project by the author,
TypeScript/LangChain/Cucumber.js); components of the Citrate Network's
consensus layer; and the Polyp Framework, a universal project scaffold
for agentic IDEs. These used different agent platforms (Claude Code,
Windsurf/Cascade, Cursor, Replit Agent) and different stacks, giving
variation, but not controlled experimental conditions.

\textbf{Implementation status.} This paper describes the
\emph{builder's} methodology, not a change to the Citrate protocol. We
use \textbf{{[}Observation{]}} for patterns noticed across projects,
\textbf{{[}Practice{]}} for recommended workflows, and
\textbf{{[}Hypothesis{]}} for testable predictions. New in v3:
\textbf{{[}Codified{]}} marks practices that were subsequently made
non-negotiable in the Agentile framework (Section 4.3).

\subsection{2. Background: BDD and
Gherkin}\label{background-bdd-and-gherkin}

\textbf{2.1 Origins.} BDD was introduced by Dan North (2003--2006)
{[}1{]} to resolve recurring confusion in Test-Driven Development: where
to start, what to test, what not to test. By replacing ``test'' with
``behavior,'' North reframed the practice from verification to
specification, the key insight. BDD draws on TDD's red-green-refactor
cycle {[}3, 4{]}, Evans' Domain-Driven Design and its ubiquitous
language {[}5{]}, and acceptance-test-driven development where
acceptance criteria are the definition of done {[}6{]}.

\textbf{2.2 The Gherkin language.} Gherkin is a structured
natural-language format for behavior {[}2{]}. Its Given/When/Then
template decomposes behavior into preconditions, actions, and expected
outcomes:

\begin{verbatim}
Feature: User Authentication
  As a healthcare provider
  I want to authenticate with MFA
  So that patient data remains HIPAA-compliant

  Scenario: Successful MFA login
    Given a registered provider with email "dr@clinic.org"
    And MFA is enabled for their account
    When they submit valid credentials
    And they enter the correct MFA code
    Then they should receive an authenticated session
    And the session should expire after 30 minutes of inactivity
\end{verbatim}

Each scenario is executable, tools like Cucumber.js {[}7{]} or
pytest-bdd {[}8{]} map the steps to functions, so the feature file is
simultaneously specification, test, and documentation. That triple
function is what makes Gherkin suited to agentic development, where the
agent needs all three in one parseable format.

\textbf{2.3 The red-green-refactor cycle.} Beck's TDD cycle {[}3{]}:
\textbf{Red}, write a failing test (the behavior does not yet exist);
\textbf{Green}, write the minimum code to pass it; \textbf{Refactor},
restructure without changing behavior (the still-passing test confirms
it). The property that matters for agentic work is that each phase is a
\textbf{verifiable checkpoint}: an agent can be stopped and resumed at
any phase without losing progress, because each phase's output is
committed and independently checkable.

\subsection{3. The Agent Discipline
Problem}\label{the-agent-discipline-problem}

\textbf{3.1 Observed failure modes.} \textbf{{[}Observation{]}} Across
three projects on multiple platforms, five recurring failure modes
appeared when agents ran without BDD constraints. \textbf{Scope drift:}
given ``implement authentication,'' agents expanded scope unasked,
adding reset flows, admin panels, each addition reasonable, collectively
a much larger system than specified. \textbf{Architectural amnesia:} in
sessions over \textasciitilde2 hours, agents forgot earlier decisions, a
module specified for dependency injection would be reimplemented with
hardcoded dependencies fifty messages later, as earlier decisions fell
out of the attention window. \textbf{Silent regressions:} modifying
existing code to add features broke working functionality without
detection or report. \textbf{Stub proliferation:} agents created
placeholder stubs and lost track of them; a post-build audit of one
project found 23 forgotten stub functions. \textbf{Platform
inconsistency:} the same prose spec produced architecturally
incompatible implementations across Claude Code, Windsurf, and Cursor.

\textbf{3.2 Why BDD addresses these.} Each failure maps to a BDD
mechanism.

{\def\LTcaptype{none} % do not increment counter
\begin{longtable}[]{@{}
  >{\raggedright\arraybackslash}p{(\linewidth - 4\tabcolsep) * \real{0.3415}}
  >{\raggedright\arraybackslash}p{(\linewidth - 4\tabcolsep) * \real{0.3902}}
  >{\raggedright\arraybackslash}p{(\linewidth - 4\tabcolsep) * \real{0.2683}}@{}}
\toprule\noalign{}
\begin{minipage}[b]{\linewidth}\raggedright
Failure mode
\end{minipage} & \begin{minipage}[b]{\linewidth}\raggedright
BDD mitigation
\end{minipage} & \begin{minipage}[b]{\linewidth}\raggedright
Mechanism
\end{minipage} \\
\midrule\noalign{}
\endhead
\bottomrule\noalign{}
\endlastfoot
Scope drift & Feature file defines exact scope & Agent can only
implement specified scenarios \\
Architectural amnesia & \texttt{.feature} files persist across sessions
& Specification survives context-window limits \\
Silent regressions & Prior scenarios stay executable & New code must
pass all existing tests \\
Stub proliferation & Red phase catches unimplemented stubs & Failing
tests explicitly surface gaps \\
Platform inconsistency & Gherkin is platform-agnostic & Same
\texttt{.feature} file works across all agents \\
\end{longtable}
}

The mitigation is \textbf{structural, not behavioral}: it does not
require the agent to be ``more careful.'' The \texttt{.feature} file is
an external artifact that constrains behavior regardless of the agent's
internal state. The discipline comes from the process, not the
practitioner.

\subsection{4. The BDD-First Agentic
Workflow}\label{the-bdd-first-agentic-workflow}

\textbf{4.1 The workflow.} \textbf{{[}Practice{]}} Across the three
projects we converged on: \textbf{Step 1 Specification}, the human
writes Gherkin feature files (the primary creative act, deciding
\emph{what} the system does; the agent does not write these).
\textbf{Step 2 Red}, the agent writes step definitions and runs them;
all fail, confirming the harness works and the behavior is absent.
\textbf{Step 3 Green}, the agent writes the minimum code to pass,
constrained not to anticipate future scenarios. \textbf{Step 4
Refactor}, the agent cleans up while keeping tests green; the developer
reviews. \textbf{Step 5 Commit and advance}, the passing scenario is
committed with a conventional message (e.g.~``🟢 FEAT: authentication
--- MFA login passes''), the next scenario is activated, return to Step
2.

\textbf{4.2 The developer's role.} The developer shifts from code author
to \textbf{specification author and quality reviewer}: writing the
Gherkin, reviewing the implementation for unintended side effects, and
making architectural decisions when the spec admits several valid
implementations. This is not a claim about all contexts; for exploratory
prototyping or one-off scripts, BDD adds overhead that may not be
justified. The claim is specific: for production systems that must stay
correct over long agentic sessions, BDD-first reduces the Section 3
failure modes.

\textbf{4.3 From workflow to framework (new in v3).}
\textbf{{[}Codified{]}} The most honest update this revision can offer
is that the workflow above stopped being a personal habit and became the
project's standing engineering framework, \textbf{Agentile}. Its
non-negotiable rules operationalize this paper's argument: test-first (a
failing test precedes implementation), \textbf{no mocks or stubs} in
delivered code (a direct answer to stub proliferation, with a per-sprint
mock budget of zero), \textbf{close-with-proof} (a unit of work cannot
close without the evidence that its tests pass), and timestamped,
authored documents so decisions survive context loss. Two additions go
beyond the February draft: \textbf{mutation testing}, which checks that
the tests themselves actually bite by introducing deliberate faults and
confirming a test fails, closing the loophole where a green suite proves
little; and \textbf{red-green as the atomic unit of work}, so every
commit is a verifiable checkpoint. The software behind Papers I--III and
X--XI was built under this framework. That is not proof the methodology
is optimal, it is evidence that it is livable at the scale of a real
system, which is a different and more modest claim.

It is worth noting which rule answers which failure mode from Section 3,
because the framework was not assembled abstractly, it accreted one rule
per scar. The no-mocks rule and its zero budget are the direct answer to
stub proliferation: a delivered stub is a rule violation the build
refuses, so the 23 forgotten stubs of the earlier audit cannot
accumulate silently. Close-with-proof answers silent regressions: a unit
of work does not close on an agent's assertion that it works, it closes
on the evidence that the prior scenarios still pass, so a regression
surfaces as a failed gate rather than as a bug discovered weeks later.
Timestamped, authored documents answer architectural amnesia across
sessions, the decision record outlives the attention window. And
mutation testing answers the subtlest failure of all, the green suite
that proves nothing: by deliberately corrupting the code and requiring a
test to fail in response, it verifies that the tests would actually
catch a regression, closing the gap between ``the tests pass'' and ``the
tests would notice if they shouldn't.'' None of these rules asks the
agent to be more careful; each makes a class of carelessness
mechanically impossible to commit.

\subsection{5. BDD as Mentorship: connection to Paper
III}\label{bdd-as-mentorship-connection-to-paper-iii}

The BDD cycle maps to the Mentorship Protocol (Paper III). The test
suite is the \textbf{mentor}, encoding the specification of correct
behavior; the coding agent is the \textbf{mentee}, receiving guidance
(failing tests) and demonstrating learning (passing tests);
red-green-refactor is the \textbf{mentorship loop}, the mentor
identifies a gap (red), the mentee fills it (green), the mentee refines
(refactor).

{\def\LTcaptype{none} % do not increment counter
\begin{longtable}[]{@{}
  >{\raggedright\arraybackslash}p{(\linewidth - 4\tabcolsep) * \real{0.1667}}
  >{\raggedright\arraybackslash}p{(\linewidth - 4\tabcolsep) * \real{0.4359}}
  >{\raggedright\arraybackslash}p{(\linewidth - 4\tabcolsep) * \real{0.3974}}@{}}
\toprule\noalign{}
\begin{minipage}[b]{\linewidth}\raggedright
BDD concept
\end{minipage} & \begin{minipage}[b]{\linewidth}\raggedright
Mentorship Protocol (Paper III)
\end{minipage} & \begin{minipage}[b]{\linewidth}\raggedright
Citrate Network (Papers I--II)
\end{minipage} \\
\midrule\noalign{}
\endhead
\bottomrule\noalign{}
\endlastfoot
Feature file & Performance profile & Per-node capability metrics at
checkpoint \\
Failing test (Red) & Identified weakness & Belnap F/N state on an input
class \\
Passing test (Green) & Successful knowledge transfer & LoRA adapter
improves node accuracy \\
Refactor & Double-loop learning & Routing model restructures
coordination \\
Test suite & Central Oracle's knowledge base & Committed embeddings +
state vectors \\
Commit & Checkpoint & BFT finality checkpoint (interval 50 blocks) \\
\end{longtable}
}

\textbf{Analogy boundary.} This mapping is structural, not operational:
a \texttt{.feature} file is not literally a performance profile, and a
failing test is not literally a Belnap classification. The claim is that
both systems implement the same organizational pattern, targeted
guidance based on observed weakness, through different mechanisms.

\subsection{6. Gherkin for Adapter
Validation}\label{gherkin-for-adapter-validation}

\textbf{{[}Practice{]}} In the Citrate Network, LoRA adapters are
generated at checkpoints and distributed to nodes (Paper II §4.3).
Before an adapter is registered on-chain, it should pass a validation
suite, and we propose expressing that suite in Gherkin, so acceptance
criteria are readable by developers, auditors, and governance
participants who are not ML specialists:

\begin{verbatim}
Feature: LoRA Adapter Validation
  Scenario: Adapter improves target domain
    Given a node with baseline accuracy of <baseline> on <domain>
    When the adapter generated at checkpoint <cp> is applied
    Then accuracy on <domain> should be >= <baseline> + <threshold>

  Scenario: Adapter does not degrade other domains
    When the adapter is applied
    Then accuracy on non-target domains should not drop > 5%

  Scenario: Adapter passes fraud-proof verification
    Given the adapter hash committed at checkpoint <cp>
    When any validator replays generation from committed state
    Then the regenerated adapter should match the committed hash
\end{verbatim}

This is a proposed application, not an implemented feature; the step
definitions require the testing infrastructure of Paper II §8. The
registration and commitment path the third scenario checks does now
exist on-chain in \texttt{LoRAFactory.sol}
(\texttt{adapterModelCommitment} / \texttt{verifyAdapterAt}, Paper III),
so the ``regenerated adapter matches the committed hash'' criterion has
a concrete target to test against.

\subsection{7. Practitioner
Observations}\label{practitioner-observations}

Reported as practitioner observations, not experimental results, and not
claimed to generalize beyond the specific projects, platforms, and
developer described.

\textbf{7.1 Agent compliance.} \textbf{{[}Observation{]}} With
\texttt{.feature} files in the project root, every tested platform
recognized Gherkin without explicit instruction and generated matching
step definitions. Compliance was highest when the feature file was in
the initial prompt context and lowest when it existed only on disk. BDD
discipline works best when the specification is in the agent's attention
window, not merely available in the filesystem.

\textbf{7.2 Regression detection.} \textbf{{[}Observation{]}} In the
IKWE.ai project, switching to BDD-first mid-project produced a
noticeable drop in regressions: before adoption (first \textasciitilde3
weeks) the author manually caught roughly 4 unflagged regressions per
session; after (remaining \textasciitilde5 weeks) the suite caught most
automatically, with roughly 1 per session caught manually. These are
approximate counts from memory and logs, not controlled measurements,
and we do not know what fraction went undetected in either period.

\textbf{7.3 Context-window survival.} \textbf{{[}Observation{]}} Feature
files preserved architectural decisions across sessions: starting a
fresh context with the \texttt{.feature} files and existing suite was
enough to bring an agent up to speed, where prose onboarding was
error-prone and incomplete. The Polyp Framework was built to exploit
this, generating feature files as part of scaffolding so every session
begins with executable specifications.

\textbf{7.4 Cross-platform portability.} \textbf{{[}Observation{]}} The
same feature files produced compatible step definitions across Claude
Code, Windsurf, and Cursor; implementation \emph{style} differed
(class-based vs functional), but the behavioral specifications, and
therefore the acceptance criteria, were identical.

\subsection{8. Experimental Hypotheses}\label{experimental-hypotheses}

The observations suggest hypotheses to test under controlled conditions;
we have not run these experiments. \textbf{H1 --- BDD reduces
agent-introduced regressions:} assign 20 comparable feature tasks to
BDD-first vs free-form on the same platform, measure regressions caught
by a hidden suite, completeness, and time, with confidence intervals.
\textbf{H2 --- feature files improve cross-session continuity:} build
over three sessions, then start session four with feature-files+tests vs
a prose summary, measure architectural inconsistencies and time to first
productive commit under blind review. \textbf{H3 --- Gherkin
portability:} give five platforms the same ten feature files vs
equivalent prose, measure common-suite pass rate, API swappability, and
architectural similarity by dependency graph.

\subsection{9. Relationship to the Gradient Papers
Series}\label{relationship-to-the-gradient-papers-series}

Paper I is the system BDD was used to build. Paper II proposes adapter
validation at checkpoints, for which Section 6 offers Gherkin as the
specification language. Paper III provides the organizational- learning
theory; this paper provides the engineering practice, the test suite is
the operational realization of the Central Oracle's mentorship function.
Paper IX describes the Strobilation Pipeline, the jellyfish reproductive
process of staged differentiation with quality gates at each transition;
red-green-refactor is its engineering echo, each cycle producing a more
capable system through incremental differentiation, an inspirational
analogy, not a formal equivalence.

\subsection{10. Conclusion}\label{conclusion}

Behavior-Driven Development is a natural, effective discipline for AI
coding agents. Gherkin gives developers a structured way to express
intent that survives context-window limits, works across platforms, and
produces executable acceptance criteria; red-green-refactor provides
incremental checkpoints that prevent scope drift, detect regressions,
and surface stubs; and the BDD-as-mentorship mapping ties this practice
to the organizational-learning theory of Paper III. The new fact this
revision can add is that the methodology was livable enough to become
the framework the rest of the series' software was built under, evidence
of practicality, not a controlled proof of superiority. The observations
here are from a single practitioner across three projects; the Section 8
experiments describe how to test them rigorously. Until then, this is a
practitioner report, useful as guidance, honest about its evidential
limits.

\subsection{Acknowledgments}\label{acknowledgments}

Drafting and literature triage were assisted by AI systems; all
observations are the authors' own from real projects, and all mechanical
claims about the Citrate contracts were verified against the referenced
source on chain 40204. This work received no external funding.

\subsection{References}\label{references}

{[}1{]} North, D. (2006). Introducing BDD. \emph{Better Software
Magazine} / dannorth.net. {[}2{]} Wynne, M., \& Helmé, A. (2012).
\emph{The Cucumber Book: Behaviour-Driven Development for Testers and
Developers.} Pragmatic Bookshelf. {[}3{]} Beck, K. (2002).
\emph{Test-Driven Development: By Example.} Addison-Wesley. {[}4{]}
Astels, D. (2003). \emph{Test-Driven Development: A Practical Guide.}
Prentice Hall. {[}5{]} Evans, E. (2003). \emph{Domain-Driven Design.}
Addison-Wesley. {[}6{]} Adzic, G. (2009). \emph{Bridging the
Communication Gap: Specification by Example and Agile Acceptance
Testing.} Neuri. {[}7{]} Cucumber.js (2025). Cucumber for JavaScript.
github.com/cucumber/cucumber-js {[}8{]} pytest-bdd (2025). BDD library
for pytest. github.com/pytest-dev/pytest-bdd {[}9{]} Pereira, L., et
al.~(2018). Behavior-Driven Development benefits and challenges: reports
from an industrial study. \emph{EASE.} {[}10{]} Irshad, M., Britto, R.,
\& Petersen, K. (2021). Adapting BDD for large-scale software systems.
\emph{Journal of Systems and Software}, 177, 110944. {[}11{]} Binamungu,
L. P., Embury, S. M., \& Konstantinou, N. (2018). Maintaining BDD
specifications: challenges and opportunities. \emph{SANER.} {[}12{]}
Klosowski, L., Mendenhall, L. (2026). Citrate: Protocol Specification.
\emph{The Gradient Papers No.~I} (this series). {[}13{]} Klosowski, L.,
Mendenhall, L. (2026). Paraconsistent Consensus. \emph{The Gradient
Papers No.~II} (this series). {[}14{]} Klosowski, L., Mendenhall, L.
(2026). The Mentorship Protocol. \emph{The Gradient Papers No.~III}
(this series). {[}15{]} Klosowski, L., Mendenhall, L. (2026). The Medusa
Paradigm. \emph{The Gradient Papers No.~IX} (this series). {[}16{]}
Senge, P. M. (1990). \emph{The Fifth Discipline.} Doubleday. {[}17{]}
Hu, E. J., et al.~(2021). LoRA: Low-Rank Adaptation of Large Language
Models. \emph{ICLR 2022.} {[}18{]} Couto, T., et al.~(2022). On the
characterization of BDD adoption benefits: a multiple case study.
\emph{SBQS.}

\subsection{Appendix A: Cross-Paper Parameter Consistency (reconciled
against code, Aug
2026)}\label{appendix-a-cross-paper-parameter-consistency-reconciled-against-code-aug-2026}

{\def\LTcaptype{none} % do not increment counter
\begin{longtable}[]{@{}
  >{\raggedright\arraybackslash}p{(\linewidth - 4\tabcolsep) * \real{0.4231}}
  >{\raggedright\arraybackslash}p{(\linewidth - 4\tabcolsep) * \real{0.2692}}
  >{\raggedright\arraybackslash}p{(\linewidth - 4\tabcolsep) * \real{0.3077}}@{}}
\toprule\noalign{}
\begin{minipage}[b]{\linewidth}\raggedright
Parameter
\end{minipage} & \begin{minipage}[b]{\linewidth}\raggedright
Value
\end{minipage} & \begin{minipage}[b]{\linewidth}\raggedright
Source
\end{minipage} \\
\midrule\noalign{}
\endhead
\bottomrule\noalign{}
\endlastfoot
Block time & \textasciitilde2 s (testnet) & Paper I;
\texttt{testnet-config.toml:4} \\
Checkpoint interval & 50 blocks & Paper I; \texttt{checkpoint.rs:101} \\
BFT committee / quorum & 100 / 67 & Paper I;
\texttt{checkpoint.rs:88} \\
Adapter registry + commitment & \texttt{LoRAFactory.sol}
\texttt{0x6e564d22…} & Paper III (corrected from ``0x1003
precompile'') \\
Adapter rank (r) & 16 (default) & Paper II \\
Adapter regression threshold & 5\% (proposed) & This paper §6 \\
\end{longtable}
}

\begin{center}\rule{0.5\linewidth}{0.5pt}\end{center}

\emph{This paper is part of the Gradient Papers series, published by
Citrate Inc.} \emph{Correspondence: Larry@citrate.ai}
